\section{未来方向与治理挑战}
\subsection{Agent 训练与数据}
Graham 表示，当前的 Agent 训练还没充分利用“Agent 风格”的数据（agent chains、自我反馈、失败纠错），未来需要更多 agent-structured corpora 让模型学会规划与纠错。另一个方向是强化 human-in-the-loop 评估：通过人类与 Agent 的协作走查，让模型学习捕捉并纠正自己的错误。

他补充说，Agent training 应该把过程日志作为训练数据的一部分，让模型不仅看到最终 patch，还能看到 Plan、执行步骤与 Verify 结果。

\subsection{治理与安全}
\begin{knowledgebox}{未来研究核心}
重点在于：训练集成化 agent 行为、引入人类修正轨迹、构建更细粒度的评估（如过程日志是否对齐、补丁的可维护性）。
\end{knowledgebox}

\begin{warningbox}{治理与安全}
Agent 过度自动化可能带来高危决策：比如在无人监督下直接部署补丁、跨仓库调用敏感 API，必须引入审批、审计与回滚机制才能量产。
\end{warningbox}

\begin{importantbox}{方向拆解}
Graham 推荐的三条路径：1) 使用 agent-specific data 训练，2) 构建 / 扩展 SWE-bench 类型的评测，3) 引入实时 human-in-the-loop 反馈与安全机制。
\end{importantbox}

\subsection{本章小结}
\begin{itemize}
    \item 解决 agent 训练与评测的空白，是 Agent 进入成熟阶段的突破口。
    \item 安全、审批与人工干预不应该被视为附加，而是 Agent 生命周期的组成部分。
\end{itemize}

